\section{پاسخ سوال ۶}
\begin{stepbox}
\field{عملیات مورفولوژیک (گرادیان مورفولوژیک)}{
عملیات خواسته‌شده یعنی $\text{dilate}(A, W) - \text{erode}(A, W)$ معادل محاسبه‌ی \textbf{گرادیان مورفولوژیک \lr{(Morphological Gradient)}} یا استخراج مرز و لبه‌های شکل $A$ است.
با توجه به اینکه تصویر $A$ یک مثلث متساوی‌الاضلاع و عنصر ساختاری (پنجره $W$) یک دایره است:

\begin{itemize}
    \item \textbf{Dilation (اتساع):} مرزهای مثلث را به اندازه شعاع دایره $W$ به سمت بیرون گسترش می‌دهد. از آنجا که پنجره دایره‌ای است، گوشه‌های تیزِ بیرونیِ مثلثِ متسع‌شده، \textbf{گِرد (Rounded)} می‌شوند (به شکل کمان‌هایی از دایره با شعاع $W$).
    \item \textbf{Erosion (فرسایش):} مرزهای مثلث را به اندازه شعاع دایره $W$ به سمت داخل جمع می‌کند. در این حالت، از آنجا که اضلاع داخلی به موازات اضلاع اصلی کوچک می‌شوند، گوشه‌های داخلیِ مثلثِ فرسایش‌یافته همچنان \textbf{تیز} باقی می‌مانند (یک مثلث متساوی‌الاضلاع کوچک‌تر).
\end{itemize}

\textbf{نتیجه نهایی:} حاصل تفریق این دو، یک \textbf{نوار ضخیم (حاشیه)} پیرامون مرز مثلث اصلی است. لبه‌ی بیرونی این نوار دارای گوشه‌های منحنی و گِرد است، در حالی که لبه‌ی داخلی آن یک مثلث کامل با گوشه‌های کاملاً تیز می‌باشد.

\begin{tipbox}[تأثیر ابعاد پنجره \lr{(Window Size)}]
همان‌طور که در پروژه‌های قبلی تجربه کرده‌اید که سایز پنجره (مثل مقایسه سایز ۴۱ در برابر ۶۱) روی دقت حذف سایه‌ها تاثیر مستقیم دارد، در اینجا نیز شعاع دایره‌ی $W$ ضخامت نوار مرزی و میزان انحنای گوشه‌های بیرونی را تعیین می‌کند. استفاده از ساختارهای مدور برای حفظ تقارن در عملیات مورفولوژیک بسیار کارآمد است.
\end{tipbox}
}
\end{stepbox}
